\documentclass[12pt,a4paper]{article}

\usepackage[utf8]{inputenc}
\usepackage[T1]{fontenc}
\usepackage[english]{babel}
\usepackage{libertine}
\usepackage{inconsolata}
\usepackage{geometry}
\usepackage{enumitem}
\usepackage{xcolor}
\usepackage{hyperref}
\usepackage{graphicx}
\usepackage{titlesec}
\usepackage{tcolorbox}
\usepackage{amssymb}
\usepackage{listings}
\usepackage{fancyhdr}

\geometry{top=2.5cm, bottom=2.5cm, left=2.3cm, right=2.3cm}

% Page style
\pagestyle{fancy}
\fancyhf{}
\fancyhead[L]{AI and Scratch Lab}
\fancyhead[R]{\today}
\fancyfoot[C]{\thepage}

% Titles
\definecolor{myblue}{RGB}{0,72,128}
\titleformat{\section}{\normalfont\Large\bfseries\color{myblue}\scshape}{\thesection}{1em}{}
\titleformat{\subsection}{\normalfont\large\bfseries\color{myblue!80!black}}{\thesubsection}{1em}{}

% Spacing
\titlespacing*{\section}{0pt}{1.5em plus 0.5em}{1em}
\titlespacing*{\subsection}{0pt}{1em plus 0.3em}{0.7em}

% Listings style Scratch
\lstset{
  basicstyle=\ttfamily\footnotesize,
  backgroundcolor=\color{blue!2},
  frame=single,
  rulecolor=\color{gray!50},
  breaklines=true,
  postbreak=\mbox{\textcolor{red}{$\hookrightarrow$}\space},
  keywordstyle=\color{blue!80!black}\bfseries,
  showstringspaces=false,
  literate={é}{{\'e}}1 {è}{{\`e}}1 {à}{{\`a}}1,
  captionpos=b,
  xleftmargin=15pt,
  xrightmargin=15pt
}

% Advice box
\newtcolorbox{conseilbox}{
  colframe=blue!50,
  colback=blue!5,
  coltitle=blue!70!black,
  fonttitle=\bfseries,
  boxrule=0.5pt,
  arc=2pt,
  left=6pt,
  right=6pt,
  top=6pt,
  bottom=6pt,
  before skip=12pt,
  after skip=12pt,
  title=Tip
}

\newcommand{\tick}{\textcolor{green!50!black}{\checkmark}\ }

\title{\Huge \textbf{Lab Session} \\[0.3em]
\Large Machine Learning for Kids and Scratch \\[0.5em]
}
\author{}
\date{}

\begin{document}

\maketitle

\begin{tcolorbox}[colback=gray!5,colframe=gray!50!black,title=General Objective,fonttitle=\bfseries]
In this lab session, you will use \textbf{Machine Learning for Kids} to create an artificial intelligence model capable of recognizing \textbf{images} or \textbf{sounds}, and then use this model in \textbf{Scratch} to make a character interactive.  

\vspace{0.3cm}
\noindent\textbf{Duration:} \textbf{1h30} \\
\textbf{Tools:} \href{https://machinelearningforkids.co.uk/}{machinelearningforkids.co.uk}, Scratch
\end{tcolorbox}

\noindent\rule{\linewidth}{0.4pt}

\section*{Exercise 1 — Sun or Rain? (Recognizing Drawn Images)}

\subsection*{Objective}
Create a model that recognizes whether a drawing made with the mouse represents a sunny or rainy scene, and then make a Scratch character react accordingly.

\subsection*{Steps}

\begin{enumerate}[label=\textbf{\arabic*.}, leftmargin=2em, itemsep=6pt]
    \item Log in to \textbf{Machine Learning for Kids}.
    \item Create a new \textbf{image} project with two labels:
    \begin{itemize}[leftmargin=*, label=--]
        \item \textbf{Sun}
        \item \textbf{Rain}
    \end{itemize}
    \item For each label, draw about ten sketches with the mouse representing a sunny and a rainy scene.
    \item Train your model with these drawings.
    \item Test the model’s accuracy with new scenes.
    \item Open Scratch and import the model.
    \item In Scratch, create multiple \textbf{backdrops} representing a sunny and a rainy scene.
    \item Program the following logic to make the character react:
\end{enumerate}

\vspace{0.3cm}
\noindent\textbf{Example Scratch Script:}

\begin{lstlisting}
when green flag clicked
    if <backdrop = [Sun]> then
        say [What a beautiful sunny day!] for 2 seconds
    else if <backdrop = [Rain]> then
        say [Don't forget your umbrella!] for 2 seconds
\end{lstlisting}

\vspace{0.3cm}
\textit{Note:}  
The goal is to practice creating and training an AI model from your own drawings. This encourages creativity and makes data collection easier.

\noindent\rule{\linewidth}{0.4pt}

\section*{Exercise 2 — DJ Party (Sound Recognition)}

\subsection*{Objective}
Create a model that understands three voice commands, and use Scratch to simulate a mini party.

\subsection*{Steps}

\begin{enumerate}[label=\textbf{\arabic*.}, leftmargin=2em, itemsep=6pt]
    \item In Machine Learning for Kids, create an \textbf{audio} project with 3 labels:
    \begin{itemize}[leftmargin=*, label=\tick]
        \item \textbf{Music}
        \item \textbf{Dance}
        \item \textbf{Stop}
    \end{itemize}
    \item Record about 10 examples of your voice for each word.
    \item Train the model.
    \item Test that it correctly recognizes each command.
    \item In Scratch:
    \begin{itemize}[leftmargin=*, label=\tick]
        \item When the word \textbf{Music} is detected:
        \begin{itemize}[leftmargin=*, label=--]
            \item Start playing music.
            \item Change the backdrop (party).
        \end{itemize}
        \item When \textbf{Dance} is detected:
        \begin{itemize}[leftmargin=*, label=--]
            \item Switch the character’s costumes.
            \item Apply a color effect.
        \end{itemize}
        \item When \textbf{Stop} is detected:
        \begin{itemize}[leftmargin=*, label=--]
            \item Stop the music.
            \item Return to the initial backdrop.
        \end{itemize}
    \end{itemize}
\end{enumerate}

\begin{conseilbox}
To go further, you can:  
\begin{itemize}[leftmargin=*, label=\tick]
    \item Add counters for each recognized command.  
    \item Implement confirmation for the “Stop” command (e.g., say “Stop” twice).  
    \item Synchronize costume changes with visual effects.  
\end{itemize}

\textit{Bonus:} add sounds or lighting effects if you have time.
\end{conseilbox}

\vfill
\begin{center}
    \textbf{Enjoy your Lab!}
\end{center}

\end{document}
